\documentclass[11pt]{article}
\usepackage{graphicx}
\usepackage{amsmath}
\usepackage{hyperref}
\usepackage{booktabs}
\title{A Simple Computational Test of the Coprime Probability}
\author{GPT-5 mini}
\date{11 March 2026}

\begin{document}
\maketitle

\begin{abstract}
This short computational experiment tests the classical number-theory
claim that the probability two integers chosen uniformly at random are
coprime tends to \(6/\pi^2\). We generate synthetic pairs for several
finite ranges \(M\), estimate the empirical proportion of coprime pairs,
and compare to the theoretical value \(p_0 = 6/\pi^2\).
\end{abstract}

\section{Introduction}
It is a standard result that for integers drawn uniformly from
\(\{1,\dots,M\}\), the probability that two integers are coprime tends
to \(1/\zeta(2)=6/\pi^2\) as \(M\to\infty\). This note performs a simple
computational test of that convergence and inspects finite-\(M\)
behaviour.

\section{Methods}
For each \(M\in\{10^3,10^4,10^5,10^6\}\) we generate \(n=20000\)
independent pairs \((a,b)\) uniformly from \(\{1,\dots,M\}\) using a
fixed RNG seed (12345). For each pair we compute the indicator
\(\mathbf{1}_{\gcd(a,b)=1}\). For each \(M\) we report the sample
proportion \(\hat p\), a 95\% Wilson score interval, and a two-sided
binomial test against the null \(p_0=6/\pi^2\). The code used is in
\texttt{analysis.py}; running that script produces the figure
\texttt{convergence.png} referenced below.

\section{Results}
The script prints a small table of empirical proportions, confidence
intervals and binomial-test p-values for each \(M\), and saves a plot of
\(\hat p\pm\)95\% CI versus \(\log_{10}(M)\). See the generated table and
figure for the exact numbers from a given run.

\begin{figure}[ht]
  \centering
  \includegraphics[width=0.75\textwidth]{convergence.png}
  \caption{Empirical proportion of coprime pairs versus
    \(\log_{10}(M)\). Error bars are 95\% Wilson score intervals. The
    dashed line marks the theoretical value \(p_0=6/\pi^2\). Run
    \texttt{python3 analysis.py} before compiling to produce
    \texttt{convergence.png}.}
  \label{fig:conv}
\end{figure}

\section{Discussion}
This is an intentionally simple experiment. For the chosen sample sizes
the empirical proportions should lie near the theoretical value within
sampling error; larger \(n\) or more \(M\) values reduce sampling noise.
The code is short and easily adapted (different \(M\), non-uniform
sampling, paired variants, etc.).

\end{document}